\section{From spectral jump 4 to Monster symmetries: APP unfolding
\texorpdfstring{$2\to6\to54$}{2--6--54}, the Leech lattice,
\texorpdfstring{$V^\natural$}{V natural}, and the Moonshine corridor}
\label{sec:cierre-equivalencias-corredor-moonshine-20260825}

The corridor is not a sequence of numerical coincidences. It begins at the
central block of APP, preserves the sum--product orientation, branches into the
APP--Fock and APP--Witt couplings, and is coordinated anew through the
grading. The binding causal order is
\[\begin{aligned}
 B_c=9P_++5P_-
 &\longrightarrow 4\longrightarrow2\longrightarrow6\longrightarrow54\\
 &\longrightarrow
 \left\{
 \begin{array}{l}
  \text{Fock degree two and }[q]t_3=54,\\
  \text{APP--Witt alignment and }N_\alpha\cong\Lambda_{24}
 \end{array}
 \right.\\
 &\longrightarrow V^\natural
 \longrightarrow
 \bigl(\operatorname{Aut},\operatorname{ch},g\mapsto T_g\bigr)
 \longrightarrow\text{Moonshine}.
\end{aligned}\]
The two intermediate branches have distinct codomains. The common key \(54\)
is a graded invariant transported by the corridor; it is not a
linear map from a scalar into a vertex operator algebra.

\subsection{Typed hierarchy and \(2\) realizations of index \(54\)}

\begin{theorem}[Spectral Hierarchy and APP Unfolding]
\label{thm:cierre-vtm-jerarquia-app}
Let \(B_c=7I+2R=9P_++5P_-\) be the central block of APP.\@ If
\(M_\Pi,M_\Sigma\) are the defined modular compressions, then
\[
 9-5=4,
 \qquad (B_c)_{i\ne j}=2,
 \qquad \sum M_\Pi-\sum M_\Sigma=51-45=6,
\]
and the deployment over the nine positions gives
\[
 9\cdot6=459-405=54.
\]
The four integers pertain, respectively, to the spectral gap, the
local coupling, modular compression, and two-dimensional integration.
\end{theorem}

\begin{proof}
The spectral decomposition of \(B_c\) gives \(9\) and \(5\), whereas
\(7I+2R\) fixes the nondiagonal coefficient. The two aggregate sums are \(51\)
and \(45\). Each compressed entry represents nine positions, so that the
unfolding multiplies the difference by nine. None of the four integers
is obtained by retrospectively identifying its types.
\end{proof}

Once the orientation line \(o_{\Sigma\Pi}\), leaf parity, and the
triangular normalization have been fixed, the oriented coupling
\[
 \mathcal H_{\Sigma,\Pi}^{(m)}:
 M_{\mathrm{bal}}^{C_3,-}
 \longrightarrow
 \operatorname{End}\!\bigl(\mathcal F_2(\mathcal V_m)\bigr)
 \otimes o_{\Sigma\Pi}
\]
satisfies
\[
 \operatorname{Tr}_{o,\mathcal F_2}
 \left(\mathcal H_{\Sigma,\Pi}^{(m)}(\nu)\Gamma_2(g_m)\right)
 =-\frac{3m\,c(\nu)}2.
\]
For \(m=12\) and \(c(\delta)=-3\), the trace equals \(54\). On the other hand,
\[
 t_3(q)=q^{-1}\prod_{n\ge1}(1+q^n+q^{2n})^{-12}
       =q^{-1}-12+54q-76q^2-243q^3+\cdots,
\]
from which \([q]t_3=54\).

\begin{theorem}[Common Index Compatibility Cone]
\label{thm:cierre-vtm-indice-comun-54}
With the starting structures of each branch preserved, it holds that
\[
 \boxed{
 D_{\mathrm{APP}}
 =\operatorname{Tr}_{o,\mathcal F_2}
   \left(\mathcal H_{\Sigma,\Pi}^{(12)}(\delta)\Gamma_2(g_{12})\right)
 =[q]t_3
 =v_\alpha^2
 =\frac{196884}{5\cdot729+1}
 =54.}
\]
This equality constitutes a cone of evaluations toward the same integer. It does not
identify \(W_{\mathrm{APP}}\), the Fock space, the neighboring lattice, and
\(V^\natural\).
\end{theorem}

\begin{proof}
The APP census yields \(459-405=54\); the oriented trace yields
\(54\); expansion of the cyclotomic product yields \([q]t_3=54\); and
the radial vector of the neighbor satisfies \(v_\alpha^2=4^2\cdot2+11\cdot2=54\). Finally,
\(5\cdot729+1=3646\) and \(54\cdot3646=196884\). The equalities are compared only after each
object has been constructed and therefore do not numerically select the
corridor.
\end{proof}

\paragraph{Exact status of the hierarchy \(4\to2\to6\to54\) and of the lattice
identification.}
The hierarchy \(4\to2\to6\to54\), the APP--Fock trace, and the extraction
\([q]t_3\) are exact internal identities. Identification of the neighbor with
Leech uses the declared chamber, neighboring vector, and lattice
classification.
%
